\chapter{Structural Logic Tensor Networks}
\label{chap:SLTN}

\noindent
Structural Logic Tensor Networks (sLTN) extend the Logic Tensor Network (LTN) framework by introducing an explicit representation of the structural organization of data. While the original LTN framework provides a differentiable logical language for reasoning over collections of individuals, many learning problems involve additional forms of organization, such as temporal order, sequence position, spatial structure, or graph connectivity. In these settings, the relationships between positions may constitute part of the information that a model needs to reason about. sLTN addresses this limitation by making structural organization an explicit component of the logical language \cite{rinaldi2026sltn}.
The central idea of sLTN is to introduce structural dimensions as first-class elements of the language. Structural dimensions correspond to named axes of structured data, such as time steps in a temporal sequence, positions in a sequence, or nodes in a graph. Structural variables can range over these dimensions, while structural relations describe relationships between structural positions. Logical formulas can therefore refer not only to ordinary individuals, but also to the organization in which those individuals occur \cite{rinaldi2026sltn}.
The framework preserves the main principles of LTN. Logical symbols are grounded into tensors or differentiable functions, while logical connectives and quantifiers are interpreted through differentiable fuzzy semantics. Structural constraints can therefore participate in gradient-based optimization. sLTN should consequently be understood as an extension of the LTN framework to structured representations rather than as a replacement for its underlying neurosymbolic approach \cite{rinaldi2026sltn}.

\section{Motivation and Limitations of Standard LTN for Structured Data}
\noindent
The original LTN framework is primarily designed for settings in which data can be represented as collections of individuals. Each individual belongs to a logical domain and can be associated with tensor representations, predicates, functions, and other logical symbols. This formulation is suitable when the relevant relationships can be expressed directly in terms of the individuals being represented.
However, many datasets contain an additional organization that cannot be naturally described only in terms of independent individuals. Examples include temporal sequences, where observations are ordered in time; token sequences, where the position of an element is meaningful; and graphs, where connectivity between nodes contributes to the structure of the data. In such cases, structural information is not merely an implementation detail of the tensor representation, but may affect the meaning of the logical constraints themselves \cite{rinaldi2026sltn}.
For example, consider a sequence of observations indexed by time. A property holding at time $t$ may be related to the same property at a subsequent time $t'$. Expressing such a dependency requires the logical language to refer explicitly to the relationship between the two positions. Similarly, in a graph, a property associated with one node may need to be related to the properties of neighboring nodes.
Such relationships can be represented indirectly by encoding structural information through ordinary individuals and predicates. However, this does not make the structural organization itself an explicit component of the logical language. sLTN instead introduces structural dimensions and structural relations so that these forms of organization can be represented directly within logical formulas \cite{rinaldi2026sltn}.
This distinction is particularly relevant to the present thesis. By making structural organization explicit, sLTN provides a formal mechanism for investigating whether structural information can contribute to generalization when a model is evaluated on structural configurations that were not observed during training. The purpose of the framework is therefore not only to represent structured data, but also to make structural information available to differentiable logical reasoning.

\section{Structural Dimensions}
\noindent
A structural dimension is a named dimension associated with a particular organization of data. Typical examples include a temporal dimension representing time steps, a sequence dimension representing positions in an ordered sequence, or a graph dimension representing nodes. Structural dimensions provide an explicit representation of the axes along which structured data are organized \cite{rinaldi2026sltn}.
In sLTN, structural dimensions are part of the formal signature of the language. The signature extends the standard components of a many-sorted first-order language with structural dimensions and structural variables. The dimensional profile associated with a symbol specifies which structural dimensions are relevant to its tensorial representation \cite{rinaldi2026sltn}.
This establishes a connection between the logical representation and the underlying tensor representation. A tensor may contain one or more axes corresponding to structural dimensions, and these axes are explicitly named and tracked by the framework. Structural information can consequently be used when expressions are aligned, selected, masked, or aggregated.
For example, suppose that a video is represented as a sequence of frames. A temporal dimension $T$ can be associated with the sequence of frames. A tensor representing a batch of videos may therefore contain an axis corresponding to $T$. Rather than treating this axis as an anonymous tensor dimension, sLTN associates it with a named structural dimension that can be referenced by logical expressions.
The explicit naming of dimensions is important because the numerical shape of an axis does not determine its semantic meaning. An axis representing time is conceptually different from an axis representing image width, even if both are represented numerically as tensor dimensions. sLTN makes this structural meaning explicit within the language and its semantics \cite{rinaldi2026sltn}.

\section{Structural Variables}
\noindent
Structural variables are used to refer to individual positions of a structural dimension and are distinct from ordinary first-order variables \cite{rinaldi2026sltn}.
An ordinary first-order variable ranges over individuals belonging to a declared logical sort. For example, a variable $x$ of sort $\mathit{Image}$ may range over images in a dataset. A structural variable, in contrast, ranges over positions associated with a structural dimension. If $t$ is associated with a temporal dimension, it can represent a particular position in time.
This distinction allows a formula to refer simultaneously to an individual and to the structural position at which that individual is observed. For example, a property can be associated with an object $x$ at time $t$, while another property can be associated with the same object at a different time $t'$.
Structural variables can also be quantified. The language can therefore express statements over all or selected positions of a structural dimension. This is particularly useful when a logical constraint depends on relationships between multiple structural positions rather than on a single observation \cite{rinaldi2026sltn}.
The distinction between first-order and structural variables is consequently fundamental to sLTN: first-order variables describe the entities being reasoned about, whereas structural variables describe where those entities occur within an explicitly represented structure.

\section{Structural Relations}
\noindent
Structural dimensions identify positions, but structured reasoning also requires a mechanism for expressing relationships between those positions. sLTN introduces structural relations for this purpose.
A structural relation connects elements of one or more structural dimensions. Such relations may represent temporal adjacency, sequence ordering, or graph connectivity. They are interpreted as Boolean or fuzzy masks over structural dimensions and can therefore be incorporated into the differentiable semantics of the framework \cite{rinaldi2026sltn}.
A simple example is the temporal relation $\mathit{next}(t,t')$, which identifies pairs of consecutive positions in a temporal dimension. This relation can be used to express a constraint such as

$$
\forall t,t' \mid \mathit{next}(t,t') :
\left(A(x,t) \rightarrow A(x,t')\right).
$$

The formula states that whenever the relation $\mathit{next}(t,t')$ holds, the property $A$ should be preserved from one temporal position to the next. The relevant aspect is that temporal adjacency is represented directly by a structural relation and can consequently be referenced by the logical constraint \cite{rinaldi2026sltn}.
Structural relations therefore provide the mechanism through which dependencies between positions become accessible to logical reasoning. The same principle can be applied to other structures. In a graph, for example, a relation can identify neighboring nodes, while in a sequence it can represent consecutive positions or another predefined ordering.
Because structural relations may be represented by Boolean or fuzzy masks, they can be incorporated into the differentiable semantics of sLTN. Structural knowledge can therefore be combined with learned representations without requiring the structural component to be treated as an entirely separate symbolic preprocessing stage \cite{rinaldi2026sltn}.

\section{Syntax and Semantics}
\noindent
The formal language of sLTN extends a many-sorted first-order signature with structural components. The signature contains sorts, dimensions, constants, individual variables, functions, predicates, structural variables, and structural relations. Typing information additionally specifies the structural dimensions associated with symbols and the dimensions consumed or produced by functions and predicates \cite{rinaldi2026sltn}.
A central aspect of the framework is the separation between the syntactic description of a theory and its tensorial interpretation. Symbols are first declared in a signature. Formulas are then constructed and checked against this signature before concrete tensor representations are assigned through an interpretation \cite{rinaldi2026sltn}.
This separation allows the logical structure of a theory to be specified independently of the particular tensors used during learning. It also makes structural information subject to explicit consistency checks at the language level. Function and predicate applications must be compatible with their declared sorts and structural profiles, while structural annotations and quantifiers must refer to appropriate structural variables and dimensions.
The concrete syntax of sLTN supports logical operators and quantifiers in both symbolic and textual forms. In addition to ordinary first-order quantification, the language supports structural quantification, allowing formulas to quantify explicitly over positions belonging to structural dimensions \cite{rinaldi2026sltn}.
The resulting framework can therefore express formulas in which ordinary logical entities and structural positions coexist. A formula can describe a property of an individual, constrain that property across different structural positions, or condition a statement on a structural relation between those positions.

\section{Tensor Representation and Named Axes}
\noindent
A central characteristic of sLTN is the direct connection between structural information and tensor dimensions. The framework interprets terms and formulas as annotated tensors whose axes may correspond to ordinary variables, structural dimensions, and domain dimensions \cite{rinaldi2026sltn}.
Consider a collection of videos represented by a tensor. If a batch contains $32$ videos, each consisting of $8$ frames, the corresponding tensor may contain a batch axis, a temporal structural axis, and image-domain axes. The temporal axis can be explicitly associated with a structural dimension such as $T$.
Named structural axes allow expressions to be aligned according to their semantic dimensions rather than relying only on their numerical tensor positions. This becomes relevant when formulas combine expressions carrying different structural profiles.
Structural operations can therefore include alignment, broadcasting, selection, masking, and aggregation along explicitly identified axes. These operations provide a connection between the logical representation of structure and its tensorial implementation \cite{rinaldi2026sltn}.
The design also preserves compatibility with ordinary LTN representations. When no structural dimensions are associated with an expression, the additional structural machinery is absent and the corresponding interpretation reduces to the unstructured case. More generally, sLTN is designed so that the original LTN semantics is recovered as a special case when structural dimensions are not used \cite{rinaldi2026sltn}.

\section{Logical Reasoning over Structural Information}
\noindent
The introduction of structural dimensions and structural relations extends the types of constraints that can be expressed within the logical language. Standard LTN allows logical formulas to describe relationships between individuals and constrain their learned representations. sLTN additionally allows such formulas to refer explicitly to the organization of the underlying data \cite{rinaldi2026sltn}.
This makes it possible to express constraints involving temporal persistence, sequence dependencies, or relationships between connected structural positions. A structural relation acts as a logical condition identifying which combinations of positions should be considered related.
For example, let $A(x,t)$ denote a property of an individual $x$ at temporal position $t$. A persistence constraint can be expressed using a relation $\mathit{next}(t,t')$:

$$
\forall t,t' \mid \mathit{next}(t,t') :
\left(A(x,t) \rightarrow A(x,t')\right).
$$

The structural relation determines where the logical implication is applied. Consequently, the formula does not simply require $A$ to hold independently at all positions. Instead, it specifies a relationship between the truth of $A$ at two structurally connected positions \cite{rinaldi2026sltn}.
The same mechanism can represent more general structural dependencies. For instance, if a property depends on the order of elements in a sequence, the logical theory can refer to the corresponding structural positions and their relationships. In a graph setting, a relation describing adjacency can determine which nodes participate in a particular logical constraint.
Structural information can therefore be combined with learned representations within the same differentiable logical framework. Logical constraints remain subject to fuzzy semantics and gradient-based optimization, while their applicability can depend explicitly on the organization of the structured input \cite{rinaldi2026sltn}.

\section{Temporal Structures as a Case Study}
\noindent
Temporal data provide a natural example of the motivation behind sLTN because temporal order is an intrinsic part of the representation. The sLTN paper uses a video classification example in which a temporal dimension represents the sequence of frames in a video \cite{rinaldi2026sltn}.
In this setting, a video can be associated with a structural dimension $T$, whose positions correspond to individual frames. A frame can therefore be represented not only by its visual content but also by its position along the temporal dimension.
A structural relation such as $\mathit{next}(t,t')$ can identify consecutive frames. Logical formulas can then use this relation to impose constraints across adjacent positions. For example, if a property is expected to persist through consecutive frames, this dependency can be expressed explicitly using the temporal relation \cite{rinaldi2026sltn}.
The temporal example illustrates the main conceptual difference introduced by sLTN. Temporal organization is not treated solely as part of the underlying tensor or neural architecture; it becomes accessible to the logical language and can therefore participate directly in differentiable logical constraints.
This case is particularly relevant to the research question of this thesis because sequences provide a controlled setting in which structural organization can be manipulated independently of the underlying individual elements. Such a setting makes it possible to investigate whether explicitly representing structural information affects the ability of a neurosymbolic model to generalize to structural configurations that were not observed during training.

\section{Relationship between LTN and sLTN}
\noindent
sLTN should be understood as an extension of LTN rather than as a separate alternative framework. Both approaches rely on the same general neurosymbolic principle: logical expressions are grounded into differentiable tensorial representations, while logical constraints are evaluated using fuzzy semantics so that they can participate in gradient-based optimization \cite{badreddine2022logic,rinaldi2026sltn}.
The main difference concerns the representation of structure. Standard LTN provides a logical language for reasoning over individuals and their relationships, whereas sLTN introduces structural dimensions, structural variables, and structural relations as explicit components of the language \cite{rinaldi2026sltn}.
This distinction can also be understood in terms of tensor representation. In standard LTN, tensor dimensions primarily support the grounding of logical entities. In sLTN, selected dimensions can additionally be assigned an explicit structural meaning, which can then be referenced by the logical language.
An important consequence is that the original LTN setting remains available as a special case. When no structural dimensions are used, the additional structural machinery is absent and the corresponding unstructured LTN semantics is recovered \cite{rinaldi2026sltn}.
This relationship is relevant to the experimental comparison developed in the following chapter. Since sLTN extends the same general neurosymbolic framework while making structural information explicit, the comparison between LTN and sLTN can be used to investigate whether explicit structural representation produces measurable differences when the learning problem depends on structural information.

\section{Limitations and Open Issues of sLTN}
\noindent
Although sLTN introduces a formal mechanism for representing structural information, the framework is presented as an initial framework and software library rather than as a fully validated solution for structured neurosymbolic learning. The authors explicitly identify the need for further stabilization, testing, and systematic empirical validation beyond the illustrative examples presented in the paper \cite{rinaldi2026sltn}.
Several directions are identified as future work. One concerns the development of dedicated structured theories for concrete application domains, in which domain-specific objects, predicates, dimensions, and structural relations can be defined systematically. Another concerns extensions of the language with finitely recursive computational primitives while preserving differentiability and a well-defined tensor semantics \cite{rinaldi2026sltn}.
The paper also identifies the need for principled strategies for selecting or scheduling the power-mean exponent used in aggregation, with the objective of improving training stability and supporting more systematic hyperparameter selection. Further work is proposed on benchmarking the library components and comparing sLTN with alternative approaches for relational, temporal, spatial, and graph-based learning \cite{rinaldi2026sltn}.
These open issues are particularly relevant to the present thesis. The introduction of explicit structural representations provides a formal mechanism for reasoning about structured data, but the empirical conditions under which this representation is actually beneficial remain an open question. In particular, the existence of a structural language does not by itself establish that explicit structural reasoning improves generalization or reduces the amount of data required for learning.
The following chapter therefore moves from the formal framework to an empirical investigation. Rather than assuming that structural representations are always advantageous, the experimental methodology is designed to examine whether their effect depends on the degree to which the target task relies on structural information and on the availability of training data.
